\documentclass[aspectratio=169,10pt]{beamer}

% ------------------------------------------------
% PAQUETES
% ------------------------------------------------
\usepackage[utf8]{inputenc}
\usepackage[T1]{fontenc}
\usepackage[spanish]{babel}
\usepackage{lmodern}
\usepackage{graphicx}
\usepackage{booktabs}
\usepackage{array}
\usepackage{multirow}
\usepackage{ragged2e}
\usepackage{xcolor}
\usepackage{tikz}
\usepackage{amsmath}
\usepackage{siunitx}
\usepackage{hyperref}

% ------------------------------------------------
% CONFIGURACIÓN DE COLORES
% ------------------------------------------------
\definecolor{AzulInstitucional}{RGB}{0,83,145}
\definecolor{AzulClaro}{RGB}{224,239,250}
\definecolor{VerdeEnergia}{RGB}{72,130,76}
\definecolor{VerdeClaro}{RGB}{231,242,229}
\definecolor{NaranjaEnergia}{RGB}{211,122,28}
\definecolor{GrisTexto}{RGB}{65,65,65}

% ------------------------------------------------
% TEMA
% ------------------------------------------------
\usetheme{Madrid}
\usecolortheme{default}

\setbeamercolor{structure}{fg=AzulInstitucional}
\setbeamercolor{frametitle}{fg=white,bg=AzulInstitucional}
\setbeamercolor{title}{fg=AzulInstitucional}
\setbeamercolor{block title}{fg=white,bg=AzulInstitucional}
\setbeamercolor{block body}{bg=AzulClaro,fg=GrisTexto}
\setbeamercolor{alerted text}{fg=NaranjaEnergia}

\setbeamertemplate{navigation symbols}{}
\setbeamertemplate{itemize items}[circle]
\setbeamertemplate{enumerate items}[default]

\setbeamertemplate{footline}{
    \leavevmode%
    \hbox{%
    \begin{beamercolorbox}[wd=.75\paperwidth,ht=2.5ex,dp=1ex,leftskip=1em]{author in head/foot}%
    Diplomado en Energía Circular Inteligente 5.0 --- Módulo 3
    \end{beamercolorbox}%
    \begin{beamercolorbox}[wd=.25\paperwidth,ht=2.5ex,dp=1ex,center]{date in head/foot}%
    \insertframenumber/\inserttotalframenumber
    \end{beamercolorbox}}%
}

% ------------------------------------------------
% COMANDOS
% ------------------------------------------------
\newcommand{\idea}[1]{
\begin{block}{Idea clave}
#1
\end{block}
}

\newcommand{\actividad}[1]{
\begin{alertblock}{Actividad}
#1
\end{alertblock}
}

\newcommand{\pregunta}[1]{
\begin{exampleblock}{Pregunta orientadora}
#1
\end{exampleblock}
}

% ------------------------------------------------
% DATOS DEL DOCUMENTO
% ------------------------------------------------
\title{Módulo 3\\Energías renovables, almacenamiento y comunidades energéticas}
\subtitle{Diplomado en Energía Circular Inteligente 5.0}
\author{Facultad de Ingeniería\\Ingeniería de Sistemas e Ingeniería Mecatrónica}
\institute{Unicomfacauca}
\date{Popayán y Santander de Quilichao\\2026}

\begin{document}

% ==================================================
% PORTADA
% ==================================================
\begin{frame}
    \titlepage
\end{frame}

\begin{frame}{Propósito del módulo}
\justifying
Comparar soluciones de energía renovable y almacenamiento bajo criterios técnicos, económicos, ambientales y sociales, priorizando medidas de eficiencia energética antes de incorporar nueva capacidad de generación.

\vspace{0.4cm}

\idea{
El objetivo no es seleccionar la tecnología ``más popular'', sino formular soluciones técnicamente viables, seguras, mantenibles y sostenibles para las condiciones reales del Cauca.
}
\end{frame}

\begin{frame}{Distribución del módulo}
\centering
\begin{tabular}{lccc}
\toprule
\textbf{Componente} & \textbf{Teoría} & \textbf{Práctica} & \textbf{Total} \\
\midrule
Módulo 3 & 8 horas & 4 horas & 12 horas \\
\bottomrule
\end{tabular}

\vspace{0.7cm}

\begin{tabular}{p{0.18\textwidth} p{0.62\textwidth}}
\textbf{Sesión 1} & Fundamentos de renovables, solar fotovoltaica y evaluación del recurso. \\
\vspace{0.15cm}
\textbf{Sesión 2} & Almacenamiento, electrificación, microrredes y comunidades energéticas. \\
\vspace{0.15cm}
\textbf{Sesión 3} & Laboratorio de pre-dimensionamiento y diseño conceptual de una solución territorial. \\
\end{tabular}
\end{frame}

\begin{frame}{Resultados de aprendizaje}
Al finalizar el módulo, el participante podrá:

\begin{itemize}
    \item Caracterizar el recurso solar y relacionarlo con perfiles de demanda.
    \item Diferenciar generación fotovoltaica, almacenamiento, electrificación, respaldo y microrredes.
    \item Realizar un pre-dimensionamiento energético argumentado.
    \item Comparar alternativas de baterías y proponer una estrategia básica de operación.
    \item Formular una alternativa de energía renovable o comunidad energética con criterios de gobernanza.
    \item Identificar restricciones de seguridad, mantenimiento, sostenibilidad financiera y disposición final.
\end{itemize}
\end{frame}

\begin{frame}{Ruta de aprendizaje}
\centering
\begin{tikzpicture}[node distance=1.8cm, >=latex, every node/.style={align=center}]

\node[draw=AzulInstitucional, rounded corners, fill=AzulClaro, minimum width=2.7cm, minimum height=1cm] (ef) {1. Eficiencia\\energética};

\node[draw=VerdeEnergia, rounded corners, fill=VerdeClaro, right=of ef, minimum width=2.7cm, minimum height=1cm] (dem) {2. Perfil de\\demanda};

\node[draw=NaranjaEnergia, rounded corners, fill=orange!15, right=of dem, minimum width=2.7cm, minimum height=1cm] (ren) {3. Renovables y\\almacenamiento};

\node[draw=AzulInstitucional, rounded corners, fill=AzulClaro, below=of dem, minimum width=3.1cm, minimum height=1cm] (sol) {4. Solución viable\\y sostenible};

\draw[->, thick] (ef) -- (dem);
\draw[->, thick] (dem) -- (ren);
\draw[->, thick] (ren) -- (sol);
\draw[->, thick] (ef) -- (sol);

\end{tikzpicture}

\vspace{0.4cm}

\idea{
La generación renovable no reemplaza el diagnóstico energético. Primero se identifica cómo, cuándo y para qué se consume la energía.
}
\end{frame}

% ==================================================
% SESIÓN 1
% ==================================================
\section{Sesión 1}

\begin{frame}{Sesión 1: Renovables y energía solar fotovoltaica}
\begin{columns}
\column{0.55\textwidth}
\textbf{Duración:} 4 horas teóricas.

\vspace{0.3cm}

\textbf{Temas centrales:}
\begin{itemize}
    \item Transición energética y jerarquía de decisiones.
    \item Recursos renovables aplicables al territorio.
    \item Fundamentos de sistemas solares fotovoltaicos.
    \item Recurso solar, irradiación, orientación y sombras.
    \item Autoconsumo, autogeneración y generación distribuida.
\end{itemize}

\column{0.40\textwidth}
\begin{block}{Producto parcial}
Mapa inicial de oportunidad energética y definición preliminar del caso de estudio.
\end{block}
\end{columns}
\end{frame}

\begin{frame}{¿Por qué no empezar por instalar paneles?}
\begin{enumerate}
    \item Identificar el problema energético real.
    \item Construir una línea base de consumo y costos.
    \item Reducir pérdidas y usos ineficientes.
    \item Ajustar o desplazar cargas cuando sea posible.
    \item Evaluar generación renovable y almacenamiento.
    \item Definir responsables de operación, mantenimiento y financiación.
\end{enumerate}

\pregunta{
¿Tiene sentido instalar un sistema solar en una organización cuyo principal desperdicio energético ocurre fuera de las horas de sol?
}
\end{frame}

\begin{frame}{Alternativas de suministro y flexibilidad}
\begin{table}
\centering
\small
\begin{tabular}{p{0.20\textwidth} p{0.28\textwidth} p{0.35\textwidth}}
\toprule
\textbf{Alternativa} & \textbf{Aporte principal} & \textbf{Condiciones y limitaciones} \\
\midrule
Solar fotovoltaica & Generación eléctrica distribuida & Depende del recurso solar, área disponible, sombras, perfil de consumo y red. \\
\addlinespace
Biomasa y biogás & Aprovechamiento energético de residuos & Requiere suministro continuo, control de humedad, logística y gestión de emisiones. \\
\addlinespace
Baterías & Respaldo, desplazamiento de consumo y flexibilidad & Implican costos, degradación, seguridad y disposición final. \\
\addlinespace
Electrificación & Sustitución de usos térmicos o combustibles & Exige revisar potencia disponible, red, proceso y costos operativos. \\
\addlinespace
Gestión de demanda & Ajuste de horarios y potencia requerida & Depende de procesos flexibles y capacidad de operación. \\
\bottomrule
\end{tabular}
\end{table}
\end{frame}

\begin{frame}{Tecnologías renovables en contexto}
\begin{columns}
\column{0.48\textwidth}
\textbf{Solar fotovoltaica}
\begin{itemize}
    \item Convierte irradiación solar en electricidad.
    \item Es modular y escalable.
    \item Puede operar conectada a red o aislada.
    \item Resulta pertinente para cubiertas, bombeo, comercio, instituciones y vivienda rural.
\end{itemize}

\column{0.48\textwidth}
\textbf{Biomasa y biogás}
\begin{itemize}
    \item Aprovechan residuos orgánicos o agroindustriales.
    \item Pueden producir calor, electricidad o biogás.
    \item Requieren balances de disponibilidad y logística.
    \item Son relevantes para café, alimentos, lácteos y residuos orgánicos.
\end{itemize}
\end{columns}

\vspace{0.4cm}

\idea{
Una fuente renovable es apropiada cuando responde a una necesidad comprobada y puede operarse de manera segura y sostenida.
}
\end{frame}

\begin{frame}{Sistema solar fotovoltaico: componentes básicos}
\centering
\begin{tikzpicture}[node distance=1.5cm, >=latex, every node/.style={align=center,font=\small}]
\node[draw=VerdeEnergia, rounded corners, fill=VerdeClaro, minimum width=2.4cm, minimum height=0.85cm] (panel) {Módulos\\fotovoltaicos};
\node[draw=AzulInstitucional, rounded corners, fill=AzulClaro, right=of panel, minimum width=2.4cm, minimum height=0.85cm] (inv) {Inversor o\\controlador};
\node[draw=NaranjaEnergia, rounded corners, fill=orange!15, right=of inv, minimum width=2.4cm, minimum height=0.85cm] (load) {Cargas\\eléctricas};

\node[draw=GrisTexto, rounded corners, fill=gray!15, below=of inv, minimum width=2.4cm, minimum height=0.85cm] (bat) {Batería\\opcional};

\node[draw=AzulInstitucional, rounded corners, fill=AzulClaro, below=of load, minimum width=2.4cm, minimum height=0.85cm] (grid) {Red eléctrica\\opcional};

\draw[->, thick] (panel) -- (inv);
\draw[->, thick] (inv) -- (load);
\draw[<->, thick] (inv) -- (bat);
\draw[<->, thick] (inv) -- (grid);
\end{tikzpicture}

\vspace{0.5cm}

\begin{itemize}
    \item Los sistemas conectados a red suelen priorizar el autoconsumo.
    \item Las baterías no son obligatorias en todos los proyectos.
    \item La protección eléctrica, puesta a tierra, cableado y seguridad son componentes esenciales.
\end{itemize}
\end{frame}

\begin{frame}{Conceptos eléctricos esenciales}
\begin{columns}
\column{0.48\textwidth}
\begin{block}{Potencia}
La potencia representa la rapidez con que se utiliza o entrega energía.

\[
P = V \times I
\]

Donde \(P\) se expresa en watts (W), \(V\) en volts (V) e \(I\) en amperes (A).
\end{block}

\column{0.48\textwidth}
\begin{block}{Energía}
La energía acumulada durante un período depende de la potencia y el tiempo.

\[
E = P \times t
\]

La energía eléctrica suele expresarse en kilowatt-hora (kWh).
\end{block}
\end{columns}

\vspace{0.35cm}

\idea{
Un panel de 500 W no entrega necesariamente 500 W de forma continua. La potencia real depende de la irradiación, temperatura, orientación, sombras y condiciones del sistema.
}
\end{frame}

\begin{frame}{Irradiación solar y energía generada}
La energía aproximada de un sistema fotovoltaico puede estimarse de manera preliminar mediante:

\[
E_{\text{FV}} \approx P_{\text{instalada}} \times HSP \times PR
\]

\begin{table}
\centering
\small
\begin{tabular}{ll}
\toprule
\textbf{Variable} & \textbf{Significado} \\
\midrule
\(E_{\text{FV}}\) & Energía generada en un período, usualmente en kWh. \\
\(P_{\text{instalada}}\) & Potencia nominal de los módulos, en kW. \\
\(HSP\) & Horas sol pico promedio del período. \\
\(PR\) & Factor de desempeño: pérdidas por temperatura, suciedad, cables, inversor, sombras y otras. \\
\bottomrule
\end{tabular}
\end{table}

\alert{Esta ecuación es para pre-dimensionamiento; no sustituye un diseño definitivo ni una simulación detallada.}
\end{frame}

\begin{frame}{Ejemplo preliminar de generación}
\textbf{Caso:} cubierta de una institución educativa con un sistema de \(10\) kW.

\[
E_{\text{FV}} \approx 10 \text{ kW} \times 4.5 \text{ h/día} \times 0.80
\]

\[
E_{\text{FV}} \approx 36 \text{ kWh/día}
\]

\begin{itemize}
    \item En un mes de 30 días, la estimación sería aproximadamente \(1{,}080\) kWh.
    \item El resultado debe contrastarse con el perfil horario de consumo.
    \item Si la mayor parte de la demanda ocurre en la noche, la generación solar por sí sola no cubrirá ese consumo en el momento de uso.
\end{itemize}

\actividad{
Discuta qué información adicional sería necesaria antes de recomendar la compra del sistema.
}
\end{frame}

\begin{frame}{Orientación, inclinación y sombras}
\begin{columns}
\column{0.5\textwidth}
\textbf{Orientación e inclinación}
\begin{itemize}
    \item Buscan maximizar o adecuar la captación solar.
    \item Deben considerar latitud, tipo de cubierta, drenaje, viento y mantenimiento.
    \item La decisión no depende únicamente de la máxima generación anual.
\end{itemize}

\column{0.5\textwidth}
\textbf{Sombras}
\begin{itemize}
    \item Pueden provenir de árboles, edificios, postes, antenas o elementos de cubierta.
    \item Cambian con la hora y la época del año.
    \item Una sombra parcial puede afectar de forma significativa el rendimiento de un arreglo.
\end{itemize}
\end{columns}

\vspace{0.4cm}

\pregunta{
¿Qué ocurriría con la producción si la cubierta recibe sombra de árboles entre las 2:00 p.\,m. y las 4:00 p.\,m.?
}
\end{frame}

\begin{frame}{Autoconsumo y perfil de demanda}
\centering
\begin{tikzpicture}[x=1.1cm,y=0.6cm]
\draw[->] (0,0) -- (10,0) node[right] {Hora};
\draw[->] (0,0) -- (0,5.5) node[above] {Potencia};

\draw[thick,AzulInstitucional] plot[smooth] coordinates {
(0.5,1.0) (1.5,0.8) (2.5,1.0) (3.5,2.0) (4.5,3.5) (5.5,4.5)
(6.5,4.7) (7.5,4.1) (8.5,2.5) (9.5,1.5)};

\draw[thick,VerdeEnergia] plot[smooth] coordinates {
(0.5,0.0) (1.5,0.0) (2.5,0.2) (3.5,1.5) (4.5,3.2) (5.5,4.2)
(6.5,3.8) (7.5,2.3) (8.5,0.7) (9.5,0.0)};

\node[AzulInstitucional] at (7.5,5.1) {Demanda};
\node[VerdeEnergia] at (7.4,3.0) {Generación FV};
\end{tikzpicture}

\vspace{0.3cm}

\begin{itemize}
    \item El autoconsumo ocurre cuando la generación coincide con la demanda local.
    \item La coincidencia temporal es tan importante como la energía total mensual.
    \item Las medidas de gestión de carga pueden aumentar el aprovechamiento de la generación solar.
\end{itemize}
\end{frame}

\begin{frame}{Criterios para evaluar una alternativa solar}
\begin{table}
\centering
\small
\begin{tabular}{p{0.24\textwidth} p{0.61\textwidth}}
\toprule
\textbf{Criterio} & \textbf{Preguntas de evaluación} \\
\midrule
Técnico & ¿Existe área útil? ¿Hay sombras? ¿La red soporta la solución? ¿Qué perfil de demanda tiene el usuario? \\
\addlinespace
Económico & ¿Cuál es la inversión? ¿Qué ahorro se espera? ¿Qué costos de operación y reemplazo existen? \\
\addlinespace
Ambiental & ¿Qué materiales se usarán? ¿Cómo se gestionarán módulos, inversores y baterías al final de su vida útil? \\
\addlinespace
Social y territorial & ¿Quién opera? ¿Quién da mantenimiento? ¿Existe apropiación, capacidad de pago y gobernanza? \\
\addlinespace
Seguridad & ¿Qué protecciones, señalización, procedimientos y personal competente se requieren? \\
\bottomrule
\end{tabular}
\end{table}
\end{frame}

\begin{frame}{Cierre de sesión 1}
\begin{itemize}
    \item Las tecnologías renovables se evalúan después de comprender y mejorar el consumo.
    \item La solar fotovoltaica depende de recurso, área, sombras, perfil de demanda y condiciones de red.
    \item El pre-dimensionamiento es una estimación inicial; no equivale a diseño eléctrico definitivo.
    \item Una solución sostenible requiere operación, mantenimiento, seguridad y responsables claros.
\end{itemize}

\actividad{
Para la siguiente sesión, cada equipo debe seleccionar un caso preliminar: vivienda rural, bombeo, comercio, institución educativa, pequeña planta, hotel, acueducto o comunidad energética.
}
\end{frame}

% ==================================================
% SESIÓN 2
% ==================================================
\section{Sesión 2}

\begin{frame}{Sesión 2: Almacenamiento, microrredes y comunidades energéticas}
\begin{columns}
\column{0.55\textwidth}
\textbf{Duración:} 4 horas teóricas.

\vspace{0.3cm}

\textbf{Temas centrales:}
\begin{itemize}
    \item Almacenamiento electroquímico y operación de baterías.
    \item Electrificación de usos finales.
    \item Microrredes, respaldo y gestión de demanda.
    \item Biomasa y biogás: balances preliminares.
    \item Comunidades energéticas, gobernanza y sostenibilidad.
\end{itemize}

\column{0.40\textwidth}
\begin{block}{Producto parcial}
Alternativa tecnológica seleccionada y criterios iniciales de operación.
\end{block}
\end{columns}
\end{frame}

\begin{frame}{¿Por qué almacenar energía?}
Las baterías pueden cumplir varias funciones:

\begin{itemize}
    \item Respaldar cargas críticas ante interrupciones del servicio.
    \item Desplazar el uso de energía entre horas de generación y consumo.
    \item Mejorar el aprovechamiento local de generación fotovoltaica.
    \item Reducir picos de demanda en determinados casos.
    \item Permitir operación aislada o en microrredes.
    \item Dar continuidad a cargas como bombeo, refrigeración, comunicaciones o equipos esenciales.
\end{itemize}

\idea{
Una batería no ``crea'' energía: almacena parte de la energía producida o recibida y entrega menos energía útil de la que recibe debido a pérdidas.
}
\end{frame}

\begin{frame}{Variables fundamentales de una batería}
\begin{table}
\centering
\small
\begin{tabular}{p{0.26\textwidth} p{0.57\textwidth}}
\toprule
\textbf{Variable} & \textbf{Interpretación} \\
\midrule
Capacidad nominal & Cantidad de energía que puede almacenar en condiciones definidas, usualmente en kWh. \\
\addlinespace
Potencia & Capacidad de entregar o recibir energía rápidamente, expresada en kW. \\
\addlinespace
Estado de carga (SOC) & Porcentaje aproximado de energía almacenada respecto a la capacidad disponible. \\
\addlinespace
Profundidad de descarga (DoD) & Fracción de la batería que se utiliza; afecta su vida útil. \\
\addlinespace
Eficiencia de ciclo & Relación entre la energía recuperada y la energía utilizada para cargarla. \\
\addlinespace
Ciclos de vida & Número de ciclos antes de que la capacidad se reduzca a un nivel especificado. \\
\addlinespace
Degradación & Pérdida progresiva de capacidad causada por uso, temperatura, tiempo y operación. \\
\bottomrule
\end{tabular}
\end{table}
\end{frame}

\begin{frame}{Energía útil de un banco de baterías}
Una aproximación inicial puede expresarse como:

\[
E_{\text{útil}} = E_{\text{nominal}} \times DoD \times \eta
\]

Donde:

\begin{itemize}
    \item \(E_{\text{útil}}\): energía disponible para la carga.
    \item \(E_{\text{nominal}}\): energía nominal del banco.
    \item \(DoD\): profundidad máxima de descarga permitida.
    \item \(\eta\): eficiencia global del sistema.
\end{itemize}

\begin{exampleblock}{Ejemplo}
Un banco de \(10\) kWh, con \(DoD = 0.80\) y eficiencia global de \(0.90\), puede entregar aproximadamente:

\[
10 \times 0.80 \times 0.90 = 7.2 \text{ kWh}
\]
\end{exampleblock}
\end{frame}

\begin{frame}{Comparación conceptual de tecnologías de batería}
\begin{table}
\centering
\small
\begin{tabular}{p{0.22\textwidth} p{0.22\textwidth} p{0.46\textwidth}}
\toprule
\textbf{Tecnología} & \textbf{Fortalezas} & \textbf{Aspectos a verificar} \\
\midrule
Plomo-ácido & Costo inicial relativamente bajo, uso extendido & Menor profundidad de descarga útil, mayor peso, mantenimiento según tecnología, menor vida útil. \\
\addlinespace
Ion-litio & Mayor densidad energética, mayor eficiencia y vida útil en muchas aplicaciones & Costo inicial, sistema de gestión de batería, compatibilidad, seguridad y gestión al final de vida. \\
\addlinespace
Otras tecnologías & Pueden responder a necesidades específicas & Disponibilidad, madurez, soporte local, costos, garantía y condiciones de uso. \\
\bottomrule
\end{tabular}
\end{table}

\alert{La selección debe basarse en el caso de uso, no únicamente en capacidad nominal o precio de compra.}
\end{frame}

\begin{frame}{Seguridad en sistemas de almacenamiento}
\begin{itemize}
    \item Instalar sistemas de acuerdo con especificaciones técnicas, normas aplicables y supervisión competente.
    \item Evitar cortocircuitos, sobrecarga, descarga excesiva y exposición térmica inadecuada.
    \item Incorporar protecciones, fusibles, interruptores, ventilación y señalización según el sistema.
    \item Definir procedimientos de inspección, respuesta ante incidentes y mantenimiento.
    \item No intervenir bancos de baterías sin autorización, capacitación y elementos de protección adecuados.
    \item Planear el manejo y disposición final de las baterías.
\end{itemize}

\idea{
El almacenamiento agrega flexibilidad, pero también agrega responsabilidades técnicas, económicas, ambientales y de seguridad.
}
\end{frame}

\begin{frame}{Electrificación de usos finales}
\begin{columns}
\column{0.47\textwidth}
\textbf{Ejemplos de electrificación}
\begin{itemize}
    \item Sustitución de equipos térmicos ineficientes.
    \item Bombeo eléctrico eficiente.
    \item Cocción eléctrica en casos viables.
    \item Refrigeración eficiente.
    \item Movilidad eléctrica y recarga.
    \item Procesos productivos con control eléctrico.
\end{itemize}

\column{0.47\textwidth}
\textbf{Preguntas de evaluación}
\begin{itemize}
    \item ¿La red tiene capacidad disponible?
    \item ¿Cuál es la nueva demanda máxima?
    \item ¿Qué ocurre con los costos operativos?
    \item ¿El proceso conserva su calidad y productividad?
    \item ¿Qué medidas de eficiencia deben implementarse antes?
\end{itemize}
\end{columns}
\end{frame}

\begin{frame}{Microrredes: concepto básico}
\centering
\begin{tikzpicture}[node distance=1.4cm, >=latex, every node/.style={align=center,font=\small}]
\node[draw=VerdeEnergia, rounded corners, fill=VerdeClaro, minimum width=2cm, minimum height=0.8cm] (solar) {Solar FV};
\node[draw=GrisTexto, rounded corners, fill=gray!15, right=of solar, minimum width=2cm, minimum height=0.8cm] (bat) {Baterías};
\node[draw=AzulInstitucional, rounded corners, fill=AzulClaro, right=of bat, minimum width=2cm, minimum height=0.8cm] (loads) {Cargas\\críticas};

\node[draw=NaranjaEnergia, rounded corners, fill=orange!15, below=of bat, minimum width=2.2cm, minimum height=0.8cm] (control) {Control y\\medición};

\node[draw=AzulInstitucional, rounded corners, fill=AzulClaro, below=of solar, minimum width=2cm, minimum height=0.8cm] (grid) {Red o\\respaldo};

\draw[<->, thick] (solar) -- (bat);
\draw[<->, thick] (bat) -- (loads);
\draw[<->, thick] (bat) -- (control);
\draw[<->, thick] (grid) -- (control);
\draw[->, thick] (solar) -- (control);
\end{tikzpicture}

\vspace{0.45cm}

Una microrred integra generación, almacenamiento, cargas, control y, según el caso, conexión a la red principal o una fuente de respaldo.
\end{frame}

\begin{frame}{Gestión de demanda y cargas críticas}
\begin{columns}
\column{0.48\textwidth}
\textbf{Gestión de demanda}
\begin{itemize}
    \item Cambiar horarios de operación.
    \item Evitar coincidencia de cargas de alta potencia.
    \item Programar bombeo o refrigeración en horas de generación solar.
    \item Priorizar mantenimiento y eficiencia.
    \item Automatizar cargas flexibles.
\end{itemize}

\column{0.48\textwidth}
\textbf{Cargas críticas}
\begin{itemize}
    \item Iluminación de emergencia.
    \item Equipos de comunicación.
    \item Refrigeración de medicamentos o alimentos.
    \item Sistemas de bombeo prioritarios.
    \item Seguridad y control.
\end{itemize}
\end{columns}

\vspace{0.3cm}

\pregunta{
¿Todas las cargas de una organización requieren respaldo? ¿Cómo se priorizan?
}
\end{frame}

\begin{frame}{Biomasa y biogás: evaluación preliminar}
Antes de proponer una solución basada en biomasa se debe caracterizar:

\begin{itemize}
    \item Tipo de residuo disponible.
    \item Cantidad generada por día, semana o cosecha.
    \item Humedad y condiciones de almacenamiento.
    \item Poder calorífico o potencial de producción de biogás.
    \item Distancia de transporte y costos logísticos.
    \item Uso energético requerido: calor, electricidad, cocción, secado o combustible.
    \item Emisiones, olores, efluentes y subproductos.
    \item Responsable de operación y mantenimiento.
\end{itemize}

\idea{
No todo residuo es una fuente energética viable. La disponibilidad real, la estacionalidad y la logística determinan la factibilidad.
}
\end{frame}

\begin{frame}{Balance preliminar de biomasa}
\begin{block}{Preguntas básicas}
\begin{enumerate}
    \item ¿Cuántos kilogramos o toneladas de residuo se generan?
    \item ¿Con qué frecuencia se generan?
    \item ¿Qué fracción puede recuperarse realmente?
    \item ¿Cuál es la humedad del material?
    \item ¿Qué energía útil puede obtenerse?
    \item ¿Qué uso energético puede reemplazar?
\end{enumerate}
\end{block}

\begin{exampleblock}{Ejemplo conceptual}
Un residuo de café puede tener potencial energético, pero su uso dependerá de la cantidad disponible, humedad, requerimientos de secado, transporte, tecnología seleccionada y manejo seguro de subproductos.
\end{exampleblock}
\end{frame}

\begin{frame}{Comunidades energéticas}
Una comunidad energética puede entenderse como una forma de organización en la que actores del territorio participan en soluciones de energía con objetivos sociales, económicos, ambientales y de acceso.

\vspace{0.3cm}

\begin{columns}
\column{0.48\textwidth}
\textbf{Dimensión técnica}
\begin{itemize}
    \item Generación distribuida.
    \item Medición y monitoreo.
    \item Almacenamiento o respaldo.
    \item Calidad y continuidad del servicio.
    \item Mantenimiento de activos.
\end{itemize}

\column{0.48\textwidth}
\textbf{Dimensión social y organizativa}
\begin{itemize}
    \item Participación de la comunidad.
    \item Reglas de operación.
    \item Propiedad y distribución de beneficios.
    \item Formación de usuarios.
    \item Sostenibilidad financiera.
\end{itemize}
\end{columns}
\end{frame}

\begin{frame}{Gobernanza de una comunidad energética}
\begin{table}
\centering
\small
\begin{tabular}{p{0.25\textwidth} p{0.57\textwidth}}
\toprule
\textbf{Aspecto} & \textbf{Preguntas orientadoras} \\
\midrule
Actores & ¿Quiénes participan? ¿Usuarios, cooperativas, instituciones, empresa, alcaldía, universidad? \\
\addlinespace
Propiedad & ¿A quién pertenecen los activos? ¿Quién responde por reposición y mantenimiento? \\
\addlinespace
Beneficios & ¿Cómo se distribuyen los ahorros, ingresos o mejoras en continuidad del servicio? \\
\addlinespace
Operación & ¿Quién toma decisiones, registra incidencias y coordina el mantenimiento? \\
\addlinespace
Sostenibilidad & ¿Qué recursos cubren mantenimiento, reposición de equipos y formación? \\
\addlinespace
Participación & ¿Cómo se informa, consulta y vincula a las personas usuarias? \\
\bottomrule
\end{tabular}
\end{table}
\end{frame}

\begin{frame}{Caso territorial: bombeo para acueducto comunitario}
\textbf{Situación:} un acueducto comunitario requiere bombear agua y enfrenta costos de electricidad, interrupciones ocasionales y limitaciones de presupuesto.

\vspace{0.3cm}

\textbf{Alternativas posibles:}
\begin{itemize}
    \item Mejorar eficiencia de la bomba y revisar fugas antes de generar energía.
    \item Ajustar horarios de bombeo hacia horas solares, si el almacenamiento de agua lo permite.
    \item Evaluar un sistema fotovoltaico para autoconsumo.
    \item Evaluar baterías únicamente si las cargas críticas y las interrupciones lo justifican.
    \item Diseñar monitoreo de caudal, nivel de tanque, energía y estado de operación.
    \item Definir responsabilidades comunitarias para mantenimiento y sostenimiento.
\end{itemize}

\actividad{
Identifique dos beneficios y dos riesgos de incorporar un sistema solar con almacenamiento en este caso.
}
\end{frame}

\begin{frame}{Cierre de sesión 2}
\begin{itemize}
    \item El almacenamiento se justifica por una función concreta: respaldo, flexibilidad, autoconsumo o operación aislada.
    \item Las baterías deben evaluarse por energía útil, potencia, ciclos, eficiencia, degradación, seguridad y fin de vida.
    \item Las microrredes combinan componentes técnicos y reglas de operación.
    \item La biomasa requiere balances de disponibilidad, humedad, logística y uso final.
    \item Las comunidades energéticas necesitan gobernanza y sostenibilidad, además de tecnología.
\end{itemize}

\actividad{
Cada equipo debe llegar a la sesión práctica con datos o supuestos explícitos sobre su caso: consumo, demanda, horarios, área disponible, cargas críticas y restricciones.
}
\end{frame}

% ==================================================
% SESIÓN 3
% ==================================================
\section{Sesión 3}

\begin{frame}{Sesión 3: Laboratorio de pre-dimensionamiento}
\begin{columns}
\column{0.55\textwidth}
\textbf{Duración:} 4 horas prácticas.

\vspace{0.3cm}

\textbf{Objetivo:} construir una propuesta preliminar y reproducible de solución renovable, de almacenamiento o de comunidad energética.

\vspace{0.3cm}

\textbf{Modalidad:} equipos de 3 a 5 participantes.

\column{0.40\textwidth}
\begin{block}{Evidencia esperada}
Pre-dimensionamiento argumentado, modelo energético reproducible, esquema de operación y propuesta de sostenibilidad.
\end{block}
\end{columns}
\end{frame}

\begin{frame}{Reto de laboratorio}
Cada equipo desarrollará una propuesta para uno de los siguientes casos:

\begin{itemize}
    \item Vivienda rural con necesidades básicas de electricidad.
    \item Bombeo para riego o acueducto comunitario.
    \item Refrigeración de productos agrícolas o alimentos.
    \item Cubierta solar para institución educativa.
    \item Autoconsumo en comercio, hotel, restaurante o pequeña planta.
    \item Aprovechamiento de un residuo agroindustrial.
    \item Comunidad energética en etapa conceptual.
\end{itemize}

\pregunta{
¿Qué problema energético específico se resolverá y cuál es el indicador que permitirá demostrar una mejora?
}
\end{frame}

\begin{frame}{Secuencia de trabajo práctico}
\centering
\begin{enumerate}
    \item Definir alcance, usuarios, problema y restricciones.
    \item Construir o asumir una línea base energética documentada.
    \item Identificar medidas de eficiencia previas.
    \item Caracterizar la demanda: energía, potencia y horarios.
    \item Seleccionar una alternativa tecnológica pertinente.
    \item Realizar el pre-dimensionamiento.
    \item Evaluar operación, mantenimiento, seguridad y fin de vida.
    \item Preparar una propuesta de sostenibilidad y presentación técnica.
\end{enumerate}
\end{frame}

\begin{frame}{Actividad 1: Caracterización de demanda}
\textbf{Tiempo sugerido: 45 minutos}

\begin{table}
\centering
\small
\begin{tabular}{p{0.25\textwidth} p{0.14\textwidth} p{0.15\textwidth} p{0.18\textwidth} p{0.16\textwidth}}
\toprule
\textbf{Equipo o carga} & \textbf{Potencia} & \textbf{Cantidad} & \textbf{Horas/día} & \textbf{Energía/día} \\
\midrule
Iluminación & & & & \\
Bomba & & & & \\
Refrigeración & & & & \\
Computadores & & & & \\
Otros & & & & \\
\midrule
\textbf{Total} & & & & \\
\bottomrule
\end{tabular}
\end{table}

\vspace{0.3cm}

\[
E_{\text{día}} = \sum (P_i \times n_i \times t_i)
\]

\begin{itemize}
    \item Registrar supuestos y fuente de cada dato.
    \item Diferenciar cargas críticas, flexibles y no prioritarias.
\end{itemize}
\end{frame}

\begin{frame}{Actividad 2: Perfil horario de carga}
\textbf{Tiempo sugerido: 30 minutos}

\begin{itemize}
    \item Distribuir la energía diaria por intervalos horarios.
    \item Identificar las horas de mayor consumo.
    \item Señalar cargas que pueden trasladarse hacia el período solar.
    \item Identificar la demanda máxima aproximada.
    \item Determinar si el problema principal es consumo total, potencia máxima, interrupciones o falta de acceso.
\end{itemize}

\vspace{0.3cm}

\idea{
Dos organizaciones con el mismo consumo mensual pueden requerir soluciones completamente diferentes si sus perfiles horarios no coinciden.
}
\end{frame}

\begin{frame}{Actividad 3: Pre-dimensionamiento fotovoltaico}
\textbf{Tiempo sugerido: 45 minutos}

Para una primera aproximación:

\[
P_{\text{FV}} \approx \frac{E_{\text{objetivo}}}{HSP \times PR}
\]

\begin{itemize}
    \item Definir \(E_{\text{objetivo}}\): energía que se pretende cubrir con la solución.
    \item Seleccionar un valor documentado o asumido de horas sol pico.
    \item Utilizar un factor de desempeño conservador.
    \item Comparar la generación estimada con el perfil horario de demanda.
    \item Verificar de manera preliminar disponibilidad de área y posibles sombras.
\end{itemize}

\begin{exampleblock}{Ejemplo}
Para cubrir \(30\) kWh/día, con \(HSP=4.5\) y \(PR=0.80\):

\[
P_{\text{FV}} \approx \frac{30}{4.5 \times 0.80} = 8.33 \text{ kW}
\]
\end{exampleblock}
\end{frame}

\begin{frame}{Actividad 4: Dimensionamiento preliminar de batería}
\textbf{Tiempo sugerido: 30 minutos}

Cuando el caso requiera respaldo o desplazamiento de energía:

\[
E_{\text{batería nominal}} \approx \frac{E_{\text{respaldo}}}{DoD \times \eta}
\]

\begin{itemize}
    \item Definir qué cargas se respaldarán.
    \item Definir durante cuánto tiempo se requiere respaldo.
    \item Estimar la energía de respaldo requerida.
    \item Considerar profundidad de descarga y eficiencia.
    \item Verificar que la potencia de la batería o inversor soporte la demanda máxima de las cargas críticas.
\end{itemize}

\alert{No añadir batería automáticamente. Justificar su función y comparar su costo con alternativas de gestión de demanda o almacenamiento no eléctrico, como tanques de agua.}
\end{frame}

\begin{frame}{Actividad 5: Selección multicriterio}
\textbf{Tiempo sugerido: 30 minutos}

\begin{table}
\centering
\small
\begin{tabular}{p{0.25\textwidth} c c c c}
\toprule
\textbf{Criterio} & \textbf{Peso} & \textbf{Alternativa A} & \textbf{Alternativa B} & \textbf{Alternativa C} \\
\midrule
Costo inicial & & & & \\
Ahorro o beneficio esperado & & & & \\
Confiabilidad & & & & \\
Facilidad de mantenimiento & & & & \\
Seguridad & & & & \\
Impacto ambiental & & & & \\
Aceptación y gobernanza & & & & \\
Disponibilidad de repuestos & & & & \\
\bottomrule
\end{tabular}
\end{table}

\begin{itemize}
    \item Definir escalas de calificación.
    \item Explicar cada puntuación con evidencia, datos o supuestos.
    \item No usar la matriz para ocultar incertidumbres.
\end{itemize}
\end{frame}

\begin{frame}{Actividad 6: Esquema de operación}
\textbf{Tiempo sugerido: 25 minutos}

El equipo debe responder:

\begin{itemize}
    \item ¿Cómo circula la energía desde la fuente hasta las cargas?
    \item ¿Qué ocurre cuando hay baja generación solar?
    \item ¿Qué cargas se priorizan?
    \item ¿Qué variable se medirá: energía, potencia, irradiación, estado de batería, caudal o temperatura?
    \item ¿Quién realiza inspección y mantenimiento?
    \item ¿Qué falla probable debe ser detectada tempranamente?
\end{itemize}

\begin{block}{Sugerencia de herramientas}
Puede utilizarse un diagrama elaborado en PowerPoint, Draw.io, Canva, FreeCAD, Blender, QGIS, una hoja de cálculo o un esquema a mano digitalizado.
\end{block}
\end{frame}

\begin{frame}{Actividad 7: Sostenibilidad de la solución}
\textbf{Tiempo sugerido: 20 minutos}

\begin{columns}
\column{0.48\textwidth}
\textbf{Sostenibilidad técnica}
\begin{itemize}
    \item Plan de mantenimiento.
    \item Disponibilidad de repuestos.
    \item Personal responsable.
    \item Seguridad eléctrica.
    \item Monitoreo de desempeño.
\end{itemize}

\column{0.48\textwidth}
\textbf{Sostenibilidad organizativa}
\begin{itemize}
    \item Propiedad de los activos.
    \item Financiación de operación y reposición.
    \item Reglas de uso.
    \item Formación de usuarios.
    \item Participación de actores.
\end{itemize}
\end{columns}

\vspace{0.35cm}

\idea{
La solución no termina cuando se instala: debe poder mantenerse, repararse, financiarse y comprenderse durante su vida útil.
}
\end{frame}

\begin{frame}{Entregable del módulo}
Cada equipo presentará un documento o presentación breve que incluya:

\begin{enumerate}
    \item Descripción del caso, problema y restricciones.
    \item Línea base o supuestos energéticos documentados.
    \item Perfil de demanda y cargas prioritarias.
    \item Medidas de eficiencia recomendadas.
    \item Alternativa renovable, de almacenamiento o comunitaria seleccionada.
    \item Pre-dimensionamiento con cálculos reproducibles.
    \item Esquema de operación y medición.
    \item Riesgos, supuestos, mantenimiento y sostenibilidad.
    \item Recomendación final y aspectos que requieren validación profesional posterior.
\end{enumerate}
\end{frame}

\begin{frame}{Rúbrica sugerida de evaluación}
\centering
\small
\begin{tabular}{p{0.58\textwidth} c}
\toprule
\textbf{Criterio} & \textbf{Peso} \\
\midrule
Claridad del problema, alcance y contexto territorial & 15\% \\
Calidad de línea base, datos y supuestos & 15\% \\
Coherencia entre demanda, eficiencia y alternativa propuesta & 20\% \\
Pre-dimensionamiento y trazabilidad de cálculos & 20\% \\
Seguridad, operación, mantenimiento y fin de vida & 10\% \\
Viabilidad económica, social y organizativa & 10\% \\
Comunicación técnica y calidad de la sustentación & 10\% \\
\bottomrule
\end{tabular}
\end{frame}

\begin{frame}{Errores frecuentes que se deben evitar}
\begin{itemize}
    \item Dimensionar paneles sin conocer el consumo o el perfil horario.
    \item Proponer baterías sin justificar su función.
    \item Confundir potencia instalada con energía realmente generada.
    \item Asumir que toda la generación solar se aprovecha localmente.
    \item Ignorar sombras, área disponible, condiciones de cubierta y red eléctrica.
    \item No considerar mantenimiento, repuestos, capacitación y responsables.
    \item Presentar estimaciones como diseños definitivos.
    \item Omitir seguridad, disposición final y limitaciones de la propuesta.
\end{itemize}
\end{frame}

\begin{frame}{Conclusiones del módulo}
\begin{itemize}
    \item La eficiencia energética orienta la selección de renovables y almacenamiento.
    \item La energía solar fotovoltaica requiere análisis de recurso, demanda, área, sombras y condiciones de operación.
    \item El almacenamiento es una herramienta de flexibilidad, no una solución automática.
    \item Biomasa, electrificación, microrredes y comunidades energéticas deben evaluarse de manera sistémica.
    \item La tecnología debe estar acompañada de operación, mantenimiento, seguridad, gobernanza y sostenibilidad financiera.
    \item Todo pre-dimensionamiento debe ser reproducible, transparente en sus supuestos y validado antes de la implementación.
\end{itemize}
\end{frame}

\begin{frame}{Lecturas y herramientas sugeridas}
\begin{itemize}
    \item Ley 1715 de 2014 y Ley 2099 de 2021.
    \item CONPES 4075 de 2022.
    \item Decreto 2236 de 2023 y regulación vigente sobre comunidades energéticas.
    \item PAI-PROURE 2022--2030.
    \item Documentación técnica de módulos, inversores, controladores y baterías.
    \item PVsyst, SAM, RETScreen, HOMER u otras herramientas disponibles.
    \item QGIS para análisis territorial.
    \item Python y hojas de cálculo para perfiles de demanda, generación y sensibilidad.
\end{itemize}
\end{frame}

\begin{frame}
\centering
\vfill

{\LARGE \textcolor{AzulInstitucional}{Gracias}}

\vspace{0.5cm}

{\large Módulo 3: Energías renovables, almacenamiento y comunidades energéticas}

\vfill
\end{frame}

\end{document}